\section{An algebraic producer of integral-point bounds}
\label{sec:split-branches}

A framework becomes more useful when some of its global quantitative
inputs can themselves be produced from finite algebra.  The earlier
Siegel packet proves a square-completion certificate
\(D^2 f=A^2+R\), and the earlier general framework already extends it
to all powers \(D^m f=A^m+R\) with a degree gap
\cite{rungepacket,oldframework}.  We do not repeat that extension.
Here we give a different certificate, adapted to equations with several
polynomial branches.  It applies directly to mixed polynomials in two
variables and supplies both coordinate bounds, including the exceptional
fibres where a remainder polynomial vanishes.

The arithmetic principle is elementary: a nonzero integer factor of an
integer \(w\) has absolute value at most \(|w|\).  If two factors differ
by a polynomial that eventually grows faster than \(w\), an integral
solution cannot lie beyond an explicit cutoff.  This gives a producer
for a restricted, syntactically recognizable class of height bounds.
It makes no completeness claim for certificate search and does not supply
the missing general Siegel or Faltings height theorem
\cite{rungepacket,faltingspacket}.

\subsection{The finite certificate and its estimates}

For an integer polynomial \(Q(X)=\sum_{k=0}^s q_kX^k\), write
\[
 \|Q\|_1=\sum_{k=0}^s|q_k|,
 \qquad
 E_Q(H)=\sum_{k=0}^s |q_k|H^k.
\]
Trailing zero coefficients are removed before computing a degree.
The zero polynomial has the empty coefficient array; every degree
comparison below is made only after checking that the relevant
polynomials are nonzero.

\begin{definition}[A certificate from separated polynomial branches]
\label{def:split-cert}
The finite input consists of integer polynomials
\(A_0,\ldots,A_{m-1},R\), with \(m\geq2\) and \(R\ne0\), and indices
\(i,j<m\).  Form \(P=A_i-A_j\), and require
\[
 d=\deg P>r=\deg R.
\]
In particular \(P\ne0\), \(d\geq1\), and the chosen branches are
different.  Put
\begin{align}
 L&=\sum_{k<d}|p_k|,
 & C&=\|R\|_1,\label{eq:split-LC}\\
 H&=\max\{1,L+2C\},
 & K&=\max_{t<m}E_{A_t}(H)+E_R(H).
 \label{eq:split-HK}
\end{align}
All these conditions and numbers are computed from the coefficient
arrays.  Since \(R\ne0\), the integer \(C\) is at least one.
\end{definition}

\begin{lemma}[Polynomial estimates]
\label{lem:split-poly-estimates}
For \(|x|\leq H\), with \(H\geq1\), one has
\( |Q(x)|\leq E_Q(H)\).
If \(Q\) has degree \(s\geq1\), nonzero integer leading coefficient,
and \(M=|x|\geq1\), then
\begin{equation}
 M^s\leq |Q(x)|+
          \left(\sum_{k<s}|q_k|\right)M^{s-1}.
 \label{eq:split-dominance}
\end{equation}
Consequently a nonzero integer polynomial \(Q\) has no integer root
with \(|x|>\|Q\|_1\).
\end{lemma}

\begin{proof}
The triangle inequality applied to the finite evaluation sum gives
\[
 |Q(x)|\leq\sum_{k=0}^s|q_k||x|^k
              \leq\sum_{k=0}^s|q_k|H^k.
\]
For the second assertion, isolate the leading term.  Because
\(|q_s|\geq1\),
\[
 M^s\leq |q_sx^s|
 \leq |Q(x)|+\sum_{k<s}|q_k|M^k
 \leq |Q(x)|+
          \left(\sum_{k<s}|q_k|\right)M^{s-1}.
\]
The final comparison uses \(M\geq1\).
If \(Q(x)=0\), cancellation of the positive integer \(M^{s-1}\)
yields \(M\leq\sum_{k<s}|q_k|\leq\|Q\|_1\).
A nonzero constant polynomial has no root at all, so it supplies the
remaining case.  All cancellations concern positive integers; no real
root, asymptotic notation, or logarithm is needed.
\end{proof}

\begin{lemma}[Size of an integer factor]
\label{lem:split-factor}
If \(u,v,w\in\Z\), \(uv=w\), and \(w\ne0\), then
\(u\ne0\), \(v\ne0\), and \(|u|\leq|w|\).
If a finite product of integers is zero, some factor is zero.
\end{lemma}

\begin{proof}
Either zero factor would force \(w=0\).  A nonzero integer \(v\)
satisfies \(|v|\geq1\), so
\(|w|=|u||v|\geq|u|\).  For the last assertion, induct through the
finite product using the integer product-zero law.  This also yields a
bounded algorithm: scan the factor list for its first zero.  Integer
equality is decidable, so this finite search does not use an infinitary
choice principle.
\end{proof}

\begin{theorem}[Explicit bounds from separated branches]
\label{thm:split-bound}
For every certificate in Definition~\ref{def:split-cert}, every integer
solution of
\begin{equation}
 \prod_{t<m}\bigl(y-A_t(x)\bigr)=R(x)
 \label{eq:split-equation}
\end{equation}
satisfies \(|x|\leq H\) and \(|y|\leq K\), with the explicit
bounds in \eqref{eq:split-HK}.
There is a uniform primitive recursive function of the certificate
that returns a complete finite list of its integer solutions.
\end{theorem}

\begin{proof}
First prove the bound for \(x\).  Integer order is decidable; in the
case \(|x|\leq H\) there is nothing to do.  In the other case set
\(M=|x|>H\).  Since \(H\geq2C\), we have \(M>C\).
Lemma~\ref{lem:split-poly-estimates} applied to the nonzero
polynomial \(R\) shows that \(R(x)\ne0\).

Put \(u_t=y-A_t(x)\).  For a fixed \(t<m\), the product of all
other factors is an integer \(v_t\), and
\(u_tv_t=R(x)\ne0\).  Lemma~\ref{lem:split-factor} therefore gives
\( |u_t|\leq|R(x)|\).  Applying this to the two selected indices,
\begin{equation}
 |P(x)|=|u_j-u_i|
       \leq |u_j|+|u_i|\leq2|R(x)|.
 \label{eq:split-factor-upper}
\end{equation}
The polynomial estimates and \(r<d\) imply
\[
 M^d\leq |P(x)|+LM^{d-1},
 \qquad
 |R(x)|\leq CM^r\leq CM^{d-1}.
\]
Combining these with \eqref{eq:split-factor-upper},
\[
 M^d\leq (L+2C)M^{d-1}.
\]
The factor \(M^{d-1}\) is positive, so cancellation gives
\(M\leq L+2C\leq H\), contradicting the case assumption.
Thus \(|x|\leq H\).

For the \(y\)-bound, the cases \(R(x)\ne0\) and \(R(x)=0\)
are both necessary.  In the first case, the factor estimate applies
without any lower bound on \(|x|\), and in particular
\[
 |y-A_0(x)|\leq |R(x)|.
\]
Using the just-proved \(x\)-bound,
\[
 |y|\leq|A_0(x)|+|R(x)|
       \leq\max_{t<m}E_{A_t}(H)+E_R(H)=K.
\]
In the second case the finite product in
\eqref{eq:split-equation} is zero.
Lemma~\ref{lem:split-factor} gives an index \(t<m\) with
\(y=A_t(x)\).  Therefore
\( |y|\leq E_{A_t}(H)\leq K\).
This includes all exceptional fibres over roots of \(R\).

Finally, enumerate \(x=-H,-H+1,\ldots,H\) and, for each such \(x\),
enumerate \(y=-K,-K+1,\ldots,K\).  Retain a pair exactly when the
two integers in \eqref{eq:split-equation} are equal.  The two loop
lengths are \(2H+1\) and \(2K+1\), and their product is an explicit
bound on the number of tests.  Every returned pair is a solution by
its acceptance test.  Every solution occurs because of the two bounds
already proved.  Polynomial evaluation, the bounds, finite products,
and these bounded loops are primitive recursive uniformly in their
finite coefficient inputs.
\end{proof}

\begin{remark}[How to attach the certificate to another equation]
Suppose an integer polynomial \(F(X,Y)\) is given together with a
nonzero integer \(D\) and a coefficientwise checked identity
\[
 D F(X,Y)=\prod_{t<m}\bigl(Y-A_t(X)\bigr)-R(X).
\]
Then its integer zero set is exactly the set in
\eqref{eq:split-equation}: one implication substitutes \(F=0\),
and the converse cancels the nonzero integer \(D\).
Thus a coefficient identity connects the certificate to the intended
input equation.  Verifying the identity is a finite computation;
proving that every equation of interest admits such an identity would
be a separate theorem, and is not asserted here.
\end{remark}

For two branches this method is closely related to completing a
square.  With three or more branches, mixed powers of \(Y\) generally
remain, and the proof works directly with them.  This is an extension
of the certificate classes treated explicitly in the inspected packets,
not a claim of historical novelty for Runge's method.  It also has a
different purpose from the existing Mordell--Weil descent theorem,
which already produces generators and coordinates under actual
quotient-cover and contraction hypotheses \cite{mordellpacket}.
Neither input interface generates the missing hypotheses of the other.

\subsection{A complete family of mixed equations}

The preceding theorem gives a uniform height bound in a family whose
integral solutions can in fact be classified exactly.

\begin{theorem}[An exact classification for every number of branches]
\label{thm:split-classification}
For every integer \(m\geq3\), the complete integer solution set of
\begin{equation}
 \prod_{j=1}^{m}\bigl(y-x^j\bigr)=1
 \label{eq:split-family}
\end{equation}
consists of the following points:
\begin{enumerate}
 \item \((0,1)\) and \((1,2)\), for every \(m\);
 \item \((0,-1)\) and \((1,0)\), in addition, precisely when
       \(m\) is even;
 \item \((-1,0)\), in addition, precisely when
       \(m\bmod4\) is \(0\) or \(1\).
\end{enumerate}
The list is a primitive recursive function of \(m\).
\end{theorem}

\begin{proof}
Every integer factor of a product equal to \(1\) is \(1\) or \(-1\).
Indeed, its absolute value is at most one by
Lemma~\ref{lem:split-factor}, and it cannot be zero.
Since \(m\geq3\), the first and third factors are present, so
\[
 |x^3-x|
 =|(y-x)-(y-x^3)|\leq2.
\]
For \(x\geq2\), however,
\(x^3-x=x(x-1)(x+1)\geq2\cdot1\cdot3=6\).
For \(x\leq-2\), apply the same inequality to \(-x\) and use the
oddness of \(X^3-X\).  Hence \(x\in\{-1,0,1\}\).

When \(x=0\), the equation is \(y^m=1\).
Its integer solutions are \(y=1\), and also \(y=-1\) for even
\(m\): a number of absolute value at least two has a power of
absolute value at least \(2^m>1\), and zero has power zero.
When \(x=1\), exactly the same argument applied to
\((y-1)^m=1\) gives \(y=2\), and also \(y=0\) for even \(m\).

When \(x=-1\), the first two factors are \(y+1\) and \(y-1\).
They must both belong to \(\{-1,1\}\).  Since their difference is
two, the only possibility is \(y+1=1\), \(y-1=-1\), hence \(y=0\).
At this point every odd-indexed factor is \(1\), and every even-indexed
factor is \(-1\).  The product is consequently
\[
 (-1)^{\lfloor m/2\rfloor}.
\]
Write \(m=4q+s\), where \(0\leq s<4\).  Then
\(\lfloor m/2\rfloor=2q+\lfloor s/2\rfloor\), which is even
exactly for \(s=0,1\).  This proves the stated condition.

The same substitutions verify every listed point.  Thus the list is
sound and complete, and computing it needs only parity and remainder
modulo four.  The proof consists of finite integer inequalities and
bounded product identities, uniformly in \(m\).
\end{proof}

At \(m=3\), the equation is explicitly
\[
 y^3-(x+x^2+x^3)y^2+(x^3+x^4+x^5)y-x^6=1,
\]
and its only integer solutions are \((0,1)\) and \((1,2)\).
The general certificate uses the branches \(X\) and \(X^3\),
so that \(L=C=1\); it gives \(H=3\) and \(K=3^m+1\).
The direct argument in Theorem~\ref{thm:split-classification} sharpens
that coarse box to the exact list.  No assertion of smoothness,
irreducibility, or genus is used in this arithmetic classification.

\subsection{Primitive recursion and the induction audit}
\label{subsec:split-audit}

The certificate is numerical data for one fixed algorithm, rather than
a code for an arbitrary program to be evaluated.  Choose canonical
signed-integer codes and finite arrays.  Polynomial evaluation uses
Horner recursion.  Subtraction pads with zero; multiplication uses bounded
convolution; an \(m\)-fold product is a primitive recursion in \(m\).
Trailing-zero removal is a bounded scan.  Absolute coefficient sums,
weighted sums, finite maxima, and the rectangular solution scan are
bounded folds.  Thus the implementation is genuinely uniform in degree,
coefficient size, and the number of branches.

The proofs require the following quantifier-free invariants after
expanding their fixed primitive recursive symbols:
\begin{enumerate}
 \item A Horner prefix or suffix equals its corresponding finite
       evaluation sum; addition and convolution respect evaluation.
 \item A processed product has the asserted factorization.  Nonzero
       total product entails that every processed factor is nonzero,
       and every product of nonzero integers has absolute value at
       least one.
 \item Partial coefficient sums dominate partial absolute evaluations;
       for \(M\geq1\), a smaller exponent gives a no-larger power.
 \item The scan for a zero factor either stores a verified zero index
       or has checked every preceding factor to be nonzero.
 \item The solution-list accumulator contains exactly the successful
       pairs in the already processed part of the explicit box.
 \item At \((-1,0)\), the product through index \(j\) equals
       \((-1)^{\lfloor j/2\rfloor}\).
\end{enumerate}
Bounded statements about earlier indices are represented by Boolean
folds.  Their use does not enlarge the induction schema.
The large degree, coefficient, or enumeration bounds cause no
proof-theoretic difficulty: primitive recursion places no polynomial-time
restriction on an algorithm.

The supplied \texttt{SplitTargets.lean} defines named pending goals for
the factor bound, polynomial evaluation bound, dominance estimate, large-root
exclusion, the two height bounds, complete enumeration and uniform family
classification. Their definitions require integer coefficient arrays, normalized degree,
Horner evaluation, branch products, the Boolean certificate checker,
and the two bounds in \eqref{eq:split-HK}.  The checker must compute
\(A_i-A_j\) itself; trusting separately supplied difference coefficients
would leave a gap.

The displayed arguments give a concrete route to \(\ipra\): establish
the signed-integer and array equations, prove the local update lemmas,
apply quantifier-free induction to the stated invariants, and then use
intuitionistic case analysis on the decidable arithmetic guards.
Producing the corresponding \texttt{Prf PRAAx} objects remains part of
the formalization work.  Native Lean statements, Python evaluations, and
coefficient checks alone do not discharge that object-language obligation.

The accompanying reference program uses integer arithmetic only.  It
checks malformed certificates, expands the mixed cubic independently,
enumerates several complete certificate boxes, includes examples where
\(R(x)=0\), and compares the family classification for \(m=3,\ldots,12\)
on explicit small boxes.  These recorded checks validate implementation
behavior.  The global completeness claims rest on the proofs above,
not on the extent of a numerical search.
